\section{Encryption from LWE}

\subsection{A simple ring-based encryption scheme}
The following public-key encryption scheme, often referred to as LPR encryption, was originally presented in \cite{lpr10} (for power-of-two cyclotomic rings) with a full version in \cite{lpr13}. It is compact and easy to describe, yet impressively powerful, since the same idea is used for more advanced encryption schemes. The scheme is IND-CPA secure from the hardness of the Decision Ring-LWE problem \cite[Lemma 5.1]{lpr13}.

Public parameters are a ring $R = \ZZ[x]/(x^\nd + 1)$ for $n$ a power of two, a prime $q$ and an error distribution $\chi$ over $R$.

\begin{itemize}
    \item[-] \textsf{KeyGen.} Pick $a \randfrom R_q$ and (small) $s,e \leftarrow \chi$. Note here $s$ is sampled from the error distribution. Output $(\mathsf{sk}, \mathsf{pk}) = (s, a \cdot s + e) \in R_q \times R_q$.
    \item[-] \textsf{Enc.} Given an $\nd$-bit message $z \in \{ 0,1 \}^\nd$, consider it as an element of $R$ where the bits are coefficients of a polynomial. Pick (small) $r, e_1, e_2 \leftarrow \chi$. Output \[ (c_1, c_2) = (a \cdot r + e_1 \nmod q, b \cdot r + e_2 + \lfloor q/2 \rceil \cdot z \nmod q), \] where $\lfloor \cdot \rceil$ denotes rounding to the closest integer. 
    \item[-] \textsf{Dec.} Given $(c_1, c_2)$, compute \[ c_2 - c_1 \cdot s = (r \cdot e - s \cdot e_1 + e_2) + \lfloor q/2 \rceil \cdot z \nmod q. \] If the error distribution and $q$ are appropriately selected, the error terms $r \cdot e - s \cdot e_1 + e_2$ will collectively have a magnitude less than $q/4$ so that the bits of $z$ can be recovered correctly by rounding to the closest of $0$ and $\lfloor q/2 \rceil$.
\end{itemize}
LPR is also sometimes called ``Noisy ElGamal'' due to it's similarity to the ElGamal cryptosystem. The random $a$ is analogous to a generator of a cyclic multiplicative group, and exponentiation is replaced by noisy products. 

\subsection{ML-KEM}
At a high level, the now standardised Module Lattice-based Key Encapsulation Mechanism (ML-KEM, also known as CRYSTALS-Kyber, or simply Kyber) uses a tweaked version of the LPR scheme above and applies the Fujisaki-Okamoto (FO) transform to it, turning it into an IND-CCA2-secure key encapsulation mechanism. 

We present the differences compared to the LPR scheme, provide a general overview of the FO transform and show how it is applied to the modified LPR scheme. Finally, we provide some instantiation details of the actual standard.

\subsubsection{An IND-CPA scheme based on LPR}
The Kyber IND-CPA-secure encryption scheme takes the ring-based LPR scheme and translates it into an MLWE-based scheme. A high-level overview of the scheme is given in Figure~\ref{fig:cpapke}. The scheme uses (coefficient-wise) rounding functions \[ \mathsf{Compress}_q(x) = \lfloor 2/q \cdot x \rceil \mathrm{mod}^+ q, \quad \mathsf{Decompress}_q(x) = \lfloor q/2 \cdot x \rceil, \]
where, again $\lfloor \cdot \rceil$ denotes rounding to the closest integer and $\mathrm{mod}^+ q$ denotes modular reduction to the set of representatives $(-q/2, q/2] \subset \ZZ$. We present a slightly simplified variant of these functions, which are only used for LWE error correction, exactly as the rounding in the LPR scheme. The full Kyber scheme also uses (more general variants of) these to compress ciphertexts, hence the names.

\begin{figure}[h!]
\centering
\renewcommand{\arraystretch}{1.15}
\begin{tabular}{|l|l|l|}
$\mathsf{CPAPKE.KeyGen}()$ &
$\mathsf{CPAPKE.Enc}(\vv{pk}, m, r)$ &
$\mathsf{CPAPKE.Dec}(\vv{sk}, c)$ \\
\hline
\rule{0pt}{4ex}% extra space below the rule
\begin{tabular}[t]{@{}l@{}}
  Generate $\vv A \in R_q^{d \times d}$\\
Sample $\vv s \in R_q^d$ from $B_{\eta_1}$ \\
Sample $\vv e \in R_q^d$ from $B_{\eta_1}$ \\
Set $\vv{pk} = \vv A \vv s + \vv e$ and $\vv{sk} = \vv s$.
\end{tabular}
&
\begin{tabular}[t]{@{}l@{}}
  Sample $\vv r \in R_q^d$ from $B_{\eta_1}$\\
  Sample $\vv e_1 \in R_q^d$ from $B_{\eta_2}$ \\
  Let $\vv c_1 = \vv A^T \vv r + \vv e_1$ \\
  Let $c_2 = \vv{pk}^T \vv r + e_2 +  \mathsf{Decompress}_q(m)$
\end{tabular}
&
\begin{tabular}[t]{@{}l@{}}
  Let $m = \mathsf{Compress}_q(c_2 - \vv s^T \vv c_1)$
\end{tabular}
\\
return $(\vv{pk},\vv{sk})$. & return $c = (\vv c_1, c_2)$. & return $m$.
\end{tabular}
\caption{The Kyber IND-CPA-secure public-key encryption scheme.}
\label{fig:cpapke}
\end{figure}

Apart from that ring elements in $R_q$ are replaced by module elements $R_q^d$ with the corresponding changes, there are two major differences. The first is that the entries of $\vv A$ are not uniformly random and instead generated by inputting (the hash of) a 32-byte seed into an extendable-output function $\mathsf{XOF}$, an approach originally used in [NewHope], directly in the NTT domain. This mainly done for efficiency reasons as discussed in the NewHope paper.

The second is that the discrete Gaussians or rounded continuous Gaussians often used in theoretical works are replaced by two centred binomial distributions $B_{\eta_1}$ and $B_{\eta_2}$. Despite being used in early implementations, the use of Gaussian noise proved to be inefficient and/or vulnerable to timing attacks. This is why the centred binomial distribution, which is easy, efficient and secure to sample from, is used instead (cite Alg implementation paper).

\subsubsection{Turning it into a KEM: FO transform}

% Shorthands for the notation
\newcommand{\K}{\mathcal{K}}
\newcommand{\E}{\mathcal{E}}
\newcommand{\D}{\mathcal{D}}
\newcommand{\hy}{\mathrm{hy}}
\newcommand{\asym}{\mathrm{asym}}
\newcommand{\sym}{\mathrm{sym}}


\begin{figure}
\centering
\renewcommand{\arraystretch}{1.15}
\begin{tabular}{|l|l|l|}
Key-Generation $\K^{\hy}(1^k)$ &
Encryption $\E^{\hy}_{pk}(m)$ &
Decryption $\D^{\hy}_{sk}(c_1,c_2)$ \\
\hline
\rule{0pt}{4ex}% extra space below the rule
$(pk,sk) \leftarrow \K^{\asym}(1^k)$.
&
\begin{tabular}[t]{@{}l@{}}
  $\sigma \leftarrow_R \mathtt{MSPC}^{\asym}$. \\
  $r_1 := H(\sigma, m)$. \\
  $r_2 := G(\sigma)$. \\
  $c_1 := \E^{\asym}_{pk}(\sigma; r_1)$. \\
  $c_2 \leftarrow \E^{\sym}_{r_2}(m)$.
\end{tabular}
&
\begin{tabular}[t]{@{}l@{}}
  $\hat\sigma := \D^{\asym}_{sk}(c_1)$. \\
  $\hat{r_2} := G(\hat\sigma)$. \\
  $\hat m := \D^{\sym}_{\hat{r_2}}(c_2)$. \\
  $\hat{r_1} := H(\hat\sigma, \hat m)$. \\
  If $c_1 == \E^{\asym}_{pk}(\hat\sigma; \hat{r_1})$ \\
  \quad then $m := \hat m$, else $m := \perp$.
\end{tabular}
\\
return $(pk,sk)$. & return $(c_1,c_2)$. & return $m$.
\end{tabular}
\caption{Hybrid encryption scheme}
\label{fig:hybrid}
\end{figure}

\subsubsection{Instantiation details}

\subsection{Homomorphic encryption and beyond}

